% Generated from locale/tr/content/set-theory/ordinals/opps.tex; do not hand-edit.


\olfileid[tr]{sth}{ordinals}{opps}
\olsection{Ardıl ve Limit Sıra Sayıları}

Sonraki birkaç bölümde sıra sayılarını çok kullanacağız. Bu nedenle bazı
basit kavramları tanıtmak yararlı olacaktır.

\begin{defn}
Her $\alpha$ sıra sayısı için onun \emph{ardılı}
$\ordsucc{\alpha}=\alpha\cup\{\alpha\}$'dır. Bir $\beta$ sıra sayısı için
$\ordsucc{\beta}=\alpha$ ise $\alpha$'ya bir \emph{ardıl} sıra sayısı
deriz. $\alpha$ ne boş ne de ardıl bir sıra sayısıysa ona bir \emph{limit}
sıra sayısı deriz.
\end{defn}
\noindent
Aşağıdaki sonuç bunun doğru \emph{ardıl} kavramı olduğunu gösterir:

\begin{prop}
Her $\alpha$ sıra sayısı için:
\begin{enumerate}
	\item $\alpha \in \ordsucc{\alpha}$;
	\item $\ordsucc{\alpha}$ bir sıra sayısıdır;
	\item $\alpha \in \beta \in \ordsucc{\alpha}$ olan hiçbir $\beta$
	sıra sayısı yoktur.
	\end{enumerate}
\end{prop}

\begin{proof}
Açıktır ki $\alpha\in\alpha\cup\{\alpha\}=\ordsucc{\alpha}$ olur.
Aynı biçimde $\ordsucc{\alpha}$ geçişli bir sıra sayıları kümesidir ve
dolayısıyla \olref[basic]{corordtransitiveord} gereği bir sıra sayısıdır.
$\alpha\in\beta\in\ordsucc{\alpha}$ olması da olanaksızdır; çünkü bu
durumda ya $\beta\in\alpha$ ya da $\beta=\alpha$ olur ve
\olref[basic]{ordordered} ile çelişir.
\end{proof}
\noindent
Bu aynı zamanda sonlu ötesi tümevarımla ispatın bir çeşidine izin verir:

\begin{thm}[Basit Sonlu Ötesi Tümevarım]\ollabel{simpletransrecursion}
$\phi(x)$ şu koşulları sağlayan bir formül olsun:
\begin{enumerate}
	\item $\phi(\emptyset)$;
	\item her $\alpha$ sıra sayısı için $\phi(\alpha)$ ise
	$\phi(\ordsucc{\alpha})$; ve
	\item $\alpha$ bir limit sıra sayısı ve
	$(\forall\beta\in\alpha)\phi(\beta)$ ise $\phi(\alpha)$.
\end{enumerate}
O zaman $\forall\alpha\phi(\alpha)$ olur.
\end{thm}

\begin{proof}
Karşıt-tersi ispatlayacağız. $\lnot\phi$ koşulunu sağlayan bir sıra sayısı
bulunduğunu varsayın; $\gamma$ bunların en küçüğü olsun. O zaman ya
$\gamma=\emptyset$ olur; ya $\phi(\alpha)$ olan bir $\alpha$ için
$\gamma=\ordsucc{\alpha}$ olur; ya da $\gamma$ bir limit sıra sayısıdır ve
$(\forall\beta\in\gamma)\phi(\beta)$ olur.
\end{proof}
\noindent
Son bir gösterim daha sonra yararlı olacaktır:

\begin{defn}\ollabel{defsupstrict}
$X$ bir sıra sayıları kümesiyse
$\supstrict(X)=\bigcup_{\alpha\in X}\ordsucc{\alpha}$.
\end{defn}
\noindent
Burada ``lsub'', ``en küçük kesin üst sınır''ın İngilizce kısaltmasıdır.
\footnote{Bazı kitaplar bunun için ``$\text{sup}(X)$'' kullanır. Fakat
başka kitaplar ``$\text{sup}(X)$'' ifadesini en küçük \emph{kesin olmayan}
üst sınır, yani yalnızca $\bigcup X$ için kullanır. $X$'in en büyük elemanı
$\alpha$ varsa bu kavramlar ayrışır: en küçük \emph{kesin} üst sınır
$\ordsucc{\alpha}$ iken en küçük \emph{kesin olmayan} üst sınır yalnızca
$\alpha$'dır.} Aşağıdaki sonuç bunu açıklar:

\begin{prop}
$X$ bir sıra sayıları kümesiyse $\supstrict(X)$, $X$ içindeki her sıra
sayısından büyük olan en küçük sıra sayısıdır.
\end{prop}

\begin{proof}
$Y=\Setabs{\ordsucc{\alpha}}{\alpha\in X}$ olsun; böylece
$\supstrict(X)=\bigcup Y$ olur. Sıra sayıları geçişli olduğundan ve bir sıra
sayısının her elemanı bir sıra sayısı olduğundan $\supstrict(X)$ geçişli bir
sıra sayıları kümesidir; dolayısıyla
\olref[basic]{corordtransitiveord} gereği bir sıra sayısıdır.

$\alpha\in X$ ise $\ordsucc{\alpha}\in Y$ olur; böylece
$\ordsucc{\alpha}\subseteq\bigcup Y=\supstrict(X)$ ve dolayısıyla
$\alpha\in\supstrict(X)$ olur. Öyleyse $\supstrict(X)$, $X$ içindeki her
sıra sayısından kesin olarak büyüktür.

Tersine $\alpha\in\supstrict(X)$ ise bir $\beta\in X$ için
$\alpha\in\ordsucc{\beta}\in Y$ olur; böylece $\alpha\leq\beta\in X$.
Demek ki $\supstrict(X)$, $X$ üzerindeki \emph{en küçük} kesin üst sınırdır.
\end{proof}

